% === SOLUTION - UX DESIGN ===
\chapter{Solution - UX Design}

\section{General Approach}

The UX design of AIutami follows a core principle: minimize interface complexity so that participants can focus entirely on the conversation. As a real-time voice platform, every interface element is designed to be immediately understandable without prior instruction.

The interaction follows a sequential phase model---registration, session creation or joining, lobby, active discussion, conclusion, and session closure---each with a dedicated interface that displays only the actions available at that stage. This approach reduces cognitive load: users never face unnecessary choices or ambiguous options.

During the active discussion phase, the interface centers on a single primary element: the push-to-talk button. The choice of push-to-talk over an always-open microphone is deliberate: it makes the act of ``taking the floor'' an explicit gesture, aligning the digital interaction with the natural dynamics of a moderated conversation. Visual indicators display who is currently speaking, who has reserved the next turn, and the AI moderator's status.

The AI moderator is integrated into the experience as an additional participant rather than a separate control panel. It intervenes vocally when needed---prompting quieter users, redirecting off-topic exchanges, or signaling phase transitions---without requiring users to interact with it through commands or text input. Moderation happens transparently, driven by real-time analysis of the conversation's transcription and participation patterns.

\section{User Workflows}

This section describes the two main workflows that define the user experience in AIutami. The first covers the steps from account creation to the start of a moderated discussion; the second describes the interaction dynamics during and after the discussion.

\subsection{Pre-Session Flow}

\begin{figure}[H]
    \centering
    \includegraphics[width=0.9\textwidth]{workflow-pre-session.png}
    \caption{Pre-Session Workflow}
    \label{fig:workflow-pre}
\end{figure}

The pre-session workflow covers all steps from account creation to the start of a moderated discussion.

\begin{enumerate}
    \item \textbf{Registration} --- The user creates an account by providing name, surname, email, username, and password. The account is immediately active upon successful registration.
    \item \textbf{Login} --- The user authenticates with email and password, obtaining access to the platform.
    \item \textbf{Home} --- The user accesses the main view, which presents three options: create a new session, join an existing session via invitation link, or browse past sessions.
    \item \textbf{Create Session} (path A) --- The user selects a discussion context (Murder Mystery, Academic, Therapeutic, or Workplace), enters a session title, and confirms. The system creates the session and redirects the user to the lobby as host.
    \item \textbf{Join Session} (path B) --- The user pastes an invitation link received from a host. The system validates the link and adds the user to the session's lobby.
    \item \textbf{Lobby} --- All participants wait in the lobby. The host can see the participant count and, once the required number of participants has joined, starts the session.
    \item \textbf{Session Start} --- The AI moderator delivers an introductory message explaining the rules and the discussion begins.
\end{enumerate}

\subsection{In-Session Flow}

\begin{figure}[H]
    \centering
    \includegraphics[width=0.9\textwidth]{workflow-in-session.png}
    \caption{In-Session Workflow}
    \label{fig:workflow-in}
\end{figure}

The in-session workflow describes the interaction dynamics from the start of the discussion to session closure.

\begin{enumerate}
    \item \textbf{AI Introduction} --- The moderator welcomes participants, explains the interaction rules, and opens the discussion.
    \item \textbf{Discussion Phase} --- Participants engage in conversation through a turn-taking system: one speaks via push-to-talk while others can reserve the next turn with an 8-second priority window. Throughout the discussion, the AI moderator monitors participation and intervenes when needed---delivering periodic summaries, prompting inactive users, or addressing collective silence.
    \item \textbf{Conclusion Phase} --- Triggered when all participants declare readiness or when the 30-minute timer expires. The moderator delivers a vocal summary of the discussion. In the Murder Mystery context, this is followed by a voting phase where participants select the suspect they believe is guilty.
    \item \textbf{Session Closure} --- The host closes the session manually, or the system closes it automatically after 15 seconds. A final report is generated.
    \item \textbf{Session History} --- The session appears in each participant's history, where the report can be viewed and downloaded as PDF.
\end{enumerate}

\section{User Scenarios}

This section presents a series of user scenarios that illustrate the main interactions with the platform, from registration to session closure.

\subsection{Scenario 1: A New User Registers}

Luca, a new user interested in participating in moderated discussions on AIutami, decides to sign up. He accesses the homepage and selects the ``Sign up'' option.

He fills in the registration form with his first name, last name, email address, username, and password. After submitting the form, AIutami verifies that the chosen username and email address are not already associated with an existing account. Upon successful validation, Luca's account is immediately created and active.

Luca is redirected to the login page, where he enters his credentials to access the platform for the first time. Once authenticated, he reaches the home page, from which he can create a new voice session, join a session via an invitation link, or browse the history of completed sessions.

\subsection{Scenario 2: A Registered User Logs In}

Marco, an existing user of AIutami, wants to access his account. From the homepage he selects ``Log in'' and enters his email and password. The system verifies his credentials and, once authenticated, redirects him to the home page. From here he can view active sessions, create a new one, join an existing session through an invitation link shared by another participant, or browse the history of past sessions.

\subsection{Scenario 3: Creating a New Session}

Luca, who has just logged in to AIutami, wants to start a new session to discuss with other users. From the home page, he selects ``Create new session'' and is redirected to the session configuration page.

Here he chooses one of the available discussion contexts---Murder Mystery, Academic, Therapeutic, or Workplace---and enters a title for the session. Once confirmed, the system creates the session and redirects Luca to the lobby as host.

From the lobby, Luca generates an invitation link and shares it with the other participants he wants to invite. The lobby displays in real time the list of participants who have joined. Luca waits until the required number of participants has connected before starting the discussion.

\subsection{Scenario 4: Joining a Session}

Giovanni, who is already logged in to AIutami, wants to join a session created by his friend Luca. From the home page, he selects ``Join with a link'' and is redirected to the join page, where he finds a text field to paste the invitation link received from Luca. After pasting it, he confirms by clicking ``Join session.''

The system validates the link and checks that the session is still in the lobby phase and has not reached its maximum number of participants. If successful, Giovanni is redirected to the lobby, where he can see the list of participants already connected and wait for the host to start the discussion.

\subsection{Scenario 5: Starting a Session}

Marco, the host of a Murder Mystery session on AIutami, is in the lobby together with the invited participants, Luca and Giovanni. In this context, the number of participants is fixed at three, so the interface displays a counter indicating how many users have joined.

When the counter reaches the required number, Marco selects the ``Start session'' button, visible only to him as host. The session enters its active phase and all participants are redirected to the discussion interface.

The AI moderator takes the floor to introduce the session: it welcomes the participants, explains the interaction rules, and opens the discussion.

\subsection{Scenario 6: Speaking During a Session}

Marco, a participant in an active session on AIutami, wants to open the conversation. He presses the push-to-talk button to activate his microphone. The interface visually indicates that his turn is active, while the other participants can only reserve their turn to speak next.

Marco shares his thoughts and, once finished, selects the ``End turn'' button, releasing the floor. From that moment, any participant can press the push-to-talk button to take the next turn.

\subsection{Scenario 7: Reserving the Next Turn}

Giovanni, a participant in the same session, wants to speak while Marco is talking. From his interface, he presses the ``Reserve turn'' button, which is available only when another participant holds the floor.

The system registers his request. When Marco ends his turn, Giovanni receives an 8-second priority window during which only his push-to-talk button is active. Giovanni presses it within this interval and begins speaking. After he finishes and ends his turn, the floor becomes available to all participants again.

\subsection{Scenario 8: Inactive User}

Giovanni is in a session together with Luca and Marco. The discussion is about a project update and proceeds at a steady pace: Luca and Marco alternate speaking turns while Giovanni listens in silence without ever pressing the push-to-talk button.

After several minutes of inactivity, a private text notification appears on Giovanni's screen, inviting him to participate in the conversation. Giovanni reads the message but chooses not to speak and lets the others continue.

After further minutes of inactivity, the AI moderator intervenes vocally, addressing Giovanni by name and encouraging him to share his thoughts with the group. Giovanni decides to accept the invitation, presses the push-to-talk button, and contributes his perspective to the discussion.

\subsection{Scenario 9: Ready to Conclude (Murder Mystery)}

Chiara, a participant in a Murder Mystery session, has been collaborating with Matteo and Silvia to identify the culprit. After several exchanges, Chiara believes she has found a possible solution and presses the ``Ready to conclude'' button.

The system registers her action and the AI moderator announces to the group that Chiara is ready to propose a conclusion. Chiara continues participating in the discussion while Matteo and Silvia keep deliberating.

Eventually, Matteo also presses the button, and the moderator announces his declaration. When Silvia, the last remaining participant, presses ``Ready to conclude'' as well, the moderator announces that all participants are ready and the session transitions to the conclusion phase. The AI moderator delivers a vocal summary of the discussion, recapping the main points and clues that emerged.

\subsection{Scenario 10: Voting (Murder Mystery)}

After the AI moderator has delivered the discussion summary, the session enters the voting phase. Each participant sees a voting interface presenting the list of suspects. Chiara, Matteo, and Silvia must each select the character they believe is the culprit.

Chiara casts her vote first. The system registers her choice and notifies the group that she has voted, without revealing her selection. Matteo and Silvia follow. Once all three participants have voted, the system reveals the results: each participant's choice, the correct answer, and whether the group successfully identified the culprit.

\subsection{Scenario 11: Time Running Out}

During an active session on AIutami, Giulia is participating in a Workplace discussion. The maximum session duration is 30 minutes, visible in the interface.

When the elapsed time reaches 25 minutes, a text notification appears on each participant's screen, informing them that approximately five minutes remain before the conclusion phase begins. Giulia reads the message and dismisses it, returning to the discussion.

When the 30-minute mark is reached, the AI moderator intervenes vocally to announce that the discussion time has ended. The session transitions to the conclusion phase: the moderator delivers a vocal summary of the key points and outcomes of the conversation, and the session proceeds to its closing steps.

\subsection{Scenario 12: Session Closure and Report}

After all participants have cast their votes in the Murder Mystery session, Marco, as host, selects the ``Close session'' button. The system generates a final report containing the session details, voting results, and an analysis of the discussion. Marco is redirected to the session history.

Later, Sara, one of the participants, wants to share the session outcome with her professor. From the home page, she selects ``Session history'' and locates the completed session. On the detail page she can view the session creator, participants, duration, voting results, and the summary generated by the AI moderator. She downloads the report as PDF and sends it to her professor.

\section{Conversational Interaction}

The AI moderator in AIutami acts as a neutral vocal facilitator that participates in the discussion alongside human users. It does not express opinions on the topic being discussed, nor does it take sides. Its role is exclusively procedural: ensuring balanced participation, maintaining focus, and guiding the group through the session's phases.

The moderator communicates through synthesized voice messages delivered via the same audio channel used by human participants. This design choice makes the moderator's presence natural within the conversation flow, rather than relying on text overlays or separate notification systems. When the moderator speaks, the interface reflects this through the same visual indicators used for human speakers. All interventions follow a principle of minimal intrusiveness: the moderator speaks only when necessary, uses brief and indirect language, and never interrupts a participant who is speaking.

From the user's perspective, the moderator's interventions fall into distinct categories. \textbf{Session-level announcements} occur at key phase transitions: the moderator introduces the session at the start, warns participants when discussion time is running low, and delivers a closing summary when the session moves to its conclusion phase. \textbf{Periodic summaries} provide a brief vocal recap of the main points at regular intervals, optionally including a gentle redirection if the conversation has drifted or participation has become unbalanced. \textbf{Individual engagement} addresses prolonged silence from a specific participant: first through a private text notification visible only to that user, then, if the silence continues, through a vocal prompt addressing them by name. \textbf{Collective silence} handling activates when no participant takes the floor for a prolonged period, encouraging someone to restart the discussion. Finally, \textbf{contextual interventions} are triggered after speaking turns when significant issues are detected---such as monopolization, aggressive tones, or a direct request for help. Minor deviations and civil disagreements do not trigger an intervention.
